% !TEX TS-program = XeLaTeX
% !TEX encoding = UTF-8 Unicode

\chapter{系统实现与检测流程}
\label{chap:implementation}
\defaultfont

\section{系统模块划分}
基于 WiFi CSI 的脊柱姿态检测系统由数据读取模块、特征处理模块、文件级特征构建模块、模型训练模块和检测输出模块组成。各模块按照数据流顺序连接，形成从 CSI 原始数据到脊柱姿态类别输出的完整处理流程。

数据读取模块负责解析 CSI 原始文件，提取时间帧、子载波、接收天线和发射天线信息；特征处理模块负责计算接收天线相位比并进行时间相位展开；文件级特征构建模块通过滑动窗口和统计聚合生成 1620 维特征向量；模型训练模块使用带标签文件训练 ExtraTrees 分类器；检测输出模块根据训练完成的模型输出最终姿态类别。

\section{系统功能设计}
系统采用分层方式组织数据处理与分类功能。数据层负责读取原始 CSI 并形成统一的复数矩阵表示；特征层完成接收天线相位比计算、时间相位展开、窗口组织和文件级统计聚合；模型层负责 ExtraTrees 分类器的训练、保存与加载；应用层接收待检测样本并输出姿态类别。各层之间通过固定维度的数据表示连接，使训练阶段和检测阶段能够复用同一套特征处理流程。

这种功能划分将 CSI 解析、特征变换和分类决策相互分离。一方面，单个处理环节可以独立验证，便于定位数据解析或特征维度错误；另一方面，当采集设备或分类模型发生变化时，可以在保持其他环节不变的情况下替换对应模块，从而提高系统的可维护性和扩展性。

\section{训练流程}
训练阶段首先读取带标签的 CSI 文件，并按照固定随机种子在文件级划分训练集、验证集和测试集。随后，系统对每个文件执行 CSI 解析、接收天线相位比计算、相位展开、滑动窗口生成和文件级统计特征聚合，得到用于分类的 1620 维特征向量。

在模型训练阶段，ExtraTrees 分类器使用训练集文件特征学习分类规则，并通过验证集比较不同配置的表现。最终模型使用 700 棵极端随机树、类别平衡权重、平方根特征采样和固定随机种子。训练完成后，系统保存模型文件和实验指标，便于后续复测和部署。

训练完成后，系统保存分类模型，并记录数据划分、训练配置和评价指标。模型参数与实验元数据分别保存，使后续复测能够恢复与训练阶段一致的数据处理配置，同时避免测试集信息参与模型选择。

\section{检测流程}
检测阶段对新的 CSI 文件执行与训练阶段一致的数据处理流程。系统首先解析复数 CSI，然后以第 0 根接收天线为参考计算相位比，并沿时间维度展开相位；之后生成长度为 512、步长为 256 的滑动窗口，并将该文件的全部窗口聚合为 1620 维统计特征；最后加载序列化后的 ExtraTrees 模型，输出三类姿态之一。

检测流程可表示为
\begin{equation}
X_{\mathrm{raw}}\rightarrow X_{\mathrm{csi}}\rightarrow X_{\mathrm{ratio}}\rightarrow X_{\mathrm{stat}}\rightarrow \hat{y}.
\end{equation}
其中 $X_{\mathrm{ratio}}$ 表示相位比序列，$X_{\mathrm{stat}}$ 表示文件级统计特征，$\hat{y}$ 表示最终检测类别。

\section{复现性验证}
模型训练完成后，本文在独立运行环境中重新加载所保存的模型，并按照训练时记录的数据划分和特征配置对测试集进行复测。复测得到 92.04\% 的测试准确率和 92.08\% 的宏平均 F1 值，分类报告与混淆矩阵均与原始实验结果一致。这表明模型序列化过程未改变分类行为，训练与检测阶段的数据处理流程也保持一致。

\section{应用特点}
该系统不依赖摄像头图像和可穿戴设备，能够在室内 WiFi 环境下进行非接触式姿态感知。系统通过 CSI 相位比特征捕获人体姿态对无线信道的影响，并利用文件级统计特征和 ExtraTrees 集成分类器完成三类姿态检测，具有隐私友好、部署成本低和计算开销较小等特点。

